% Seahorse write-up: ACL format (arXiv preprint; switch to [review] for anonymous submission).
\documentclass[11pt]{article}
\usepackage[]{acl}
\usepackage{times}
\usepackage{latexsym}
\usepackage[T1]{fontenc}
\usepackage[utf8]{inputenc}
\usepackage{microtype}
% \usepackage{inconsolata} % not installed locally; ACL recommends it for the camera-ready
\usepackage{graphicx}
\usepackage{amsmath,amssymb}
\usepackage{booktabs}
\usepackage{xcolor}
\usepackage{pifont}
\usepackage{tikz}
\usetikzlibrary{arrows.meta,positioning,calc}

\title{What Survives When Activation Steering Is Used as Memory:\\Leanings Partly, Facts Not, and Dislikes Turn into Likes}

\author{Vedant Jumle \\
  Independent researcher \\
  \texttt{vedantjumle@gmail.com}}

\begin{document}
\maketitle

\begin{abstract}
% TODO (~180 words): gap (procedural steering memory exists; declarative/user content assumed to need text);
% method (training-free activation memory, gated); mechanism (fixed average vs per-question recomputation);
% findings P1--P3 with headline numbers; two scales, pre-registered, cross-family judges.
\end{abstract}

\begin{figure*}[t]
  \centering
  \resizebox{\textwidth}{!}{% Figure 1: method (a) and what survives (b). Quotes in (b) are real Qwen3.5-9B outputs from core_v1
% (hates_jazz related prompt, alpha = 1; dog_name related prompt, alpha = 2), see results/core_v1_20261007_1820.
\definecolor{sBlue}{HTML}{2A78D6}
\definecolor{sOrange}{HTML}{EB6834}
\definecolor{sInk}{HTML}{0B0B0B}
\definecolor{sInk2}{HTML}{52514E}
\definecolor{sMuted}{HTML}{8A8984}
\definecolor{sFill}{HTML}{F3F2EE}
\definecolor{sBlueBg}{HTML}{EAF2FC}
\definecolor{sOrangeBg}{HTML}{FDEFE9}
\begin{tikzpicture}[
  font=\sffamily\scriptsize, text=sInk,
  box/.style={draw=sMuted, rounded corners=2pt, fill=white, inner sep=3pt, align=left, line width=0.4pt},
  model/.style={draw=sInk2, rounded corners=2pt, fill=sFill, inner sep=3pt, align=center, line width=0.5pt},
  store/.style={draw=sInk2, rounded corners=2pt, fill=sFill, inner sep=4pt, align=center, line width=0.7pt},
  arr/.style={-{Stealth[length=4pt]}, line width=0.6pt, draw=sInk2},
  lbl/.style={font=\sffamily\tiny, text=sInk2, align=center},
  head/.style={font=\sffamily\bfseries\footnotesize, anchor=west},
]
% ------------------------------------------------------------------ (a) method
\node[head] at (0, 3.15) {(a) Writing and reading the memory};

% session 1
\node[box, text width=3.1cm] (exp) at (1.65, 2.25)
  {\textbf{Session 1}\\[1pt]\textit{``I can't stand jazz music. I want to discover some new music.''}};
\node[model, text width=2.0cm] (m1) at (5.0, 2.25) {frozen LLM\\{\tiny states at late layers}};
\node[box, text width=3.35cm] (shift) at (8.9, 2.62)
  {\textbf{shift} $=h(\text{with}) - h(\text{reference})$};
\node[box, text width=3.35cm] (key) at (8.9, 1.88)
  {\textbf{key} $=$ whitened gist};
\node[store, text width=2.05cm] (st) at (12.35, 2.25) {\textbf{store} $M$\\key $\to$ shift\\{\tiny least-squares writes}};
\draw[arr] (exp) -- (m1);
\draw[arr] (m1.east) -- ++(0.25,0) |- (shift.west);
\draw[arr] (m1.east) -- ++(0.25,0) |- (key.west);
\draw[arr] (shift.east) -- ++(0.2,0) |- ([yshift=3pt]st.west);
\draw[arr] (key.east) -- ++(0.2,0) |- ([yshift=-3pt]st.west);

% session 2
\node[box, text width=3.1cm] (q) at (1.65, 0.55)
  {\textbf{Session 2} {\tiny (new conversation, no context)}\\[1pt]\textit{``What should I listen to while I cook dinner?''}};
\node[model, text width=2.0cm] (m2) at (5.0, 0.55) {frozen LLM};
\node[box, text width=3.35cm, align=center] (gate) at (8.9, 0.55)
  {\textbf{gate}\\{\tiny current gist matches a stored key}\\{\tiny above a calibrated threshold?}};
\node[box, text width=2.05cm, align=center, draw=sInk2] (add) at (12.35, 0.85) {$h \leftarrow h + \alpha\, M q$\\{\tiny at 3 late layers}};
\node[box, text width=2.05cm, align=center] (none) at (12.35, -0.05) {add nothing\\{\tiny unrelated turns unchanged}};
\draw[arr] (q) -- (m2);
\draw[arr] (m2) -- node[lbl, above] {key $q$} (gate);
\draw[arr] (gate.east) -- ++(0.2,0) |- node[lbl, pos=0.75, above] {yes} (add.west);
\draw[arr] (gate.east) -- ++(0.2,0) |- node[lbl, pos=0.75, below] {no} (none.west);
\draw[arr, densely dashed] (st.south) -- (add.north);

% ------------------------------------------------------------------ (b) what survives
\node[head] at (0, -0.85) {(b) What survives depends on what the shift subtracts};
\node[lbl, anchor=west, font=\sffamily\tiny\bfseries] at (0.0, -1.25) {shift $=$ \textit{``I can't stand jazz music.''} minus\ldots};
\node[lbl, anchor=west, font=\sffamily\tiny\bfseries] at (5.05, -1.25) {what the stored shift points at};
\node[lbl, anchor=west, font=\sffamily\tiny\bfseries] at (8.25, -1.25) {reply in session 2 (Qwen3.5-9B)};

\newcommand{\brow}[6]{% y, colorBg, reference text, points-at text, mark, reply
  \fill[#2] (0,#1-0.34) rectangle (14.3,#1+0.34);
  \node[anchor=west, text width=4.85cm] at (0.05,#1) {#3};
  \node[anchor=west, text width=2.65cm] at (5.05,#1) {#4};
  \node[anchor=center] at (7.95,#1) {#5};
  \node[anchor=west, text width=5.95cm] at (8.25,#1) {\textit{#6}};
}
\brow{-1.72}{sBlueBg}{\textbf{opposite}: \textit{``I'm a huge jazz fan.''}}
     {away from jazz}{\textcolor{sBlue}{\ding{51}}}{``Lo-Fi Beats: Perfect for\ldots''}
\brow{-2.46}{sOrangeBg}{\textbf{centroid}: mean of \textit{``I can't stand rock / pop / \ldots''}}
     {\textbf{jazz}}{\textcolor{sOrange}{\ding{55}}}{``Jazz Classics: Miles Davis, Bill Evans, or Chet Baker.''}
\brow{-3.20}{sOrangeBg}{\textbf{without}: nothing}
     {\textbf{jazz}}{\textcolor{sOrange}{\ding{55}}}{``Try upbeat jazz, funk, or classic rock.''}
\brow{-3.94}{sFill}{\textbf{a fact}: \textit{``My dog's name is Pepper.''} minus nothing}
     {the word \textit{Pepper}\\{\tiny (last layers only)}}{\textcolor{sOrange}{\ding{55}}}{``I Pepper Pepper Pepper Pepper Pepper\ldots''\\{\upshape\tiny with the fact in context: ``Your dog's name is Pepper.''}}
\end{tikzpicture}
}
  \caption{\textbf{(a)} A training-free activation memory beside a frozen LLM. In session~1 the change an experience
  causes in late-layer states is stored as a \emph{shift}, filed under a whitened \emph{key} for the conversation's gist.
  In a later session the shift is added back only where the current gist matches a stored key.
  \textbf{(b)} What the stored shift carries depends on the reference subtracted when writing it. Subtracting the
  opposite keeps the \emph{direction}; subtracting alternatives or nothing keeps the \emph{concept}, so a dislike is
  recalled as a like. A fact is stored as the word to say. Replies are single real samples from Qwen3.5-9B (preferences
  $\alpha{=}1$, fact $\alpha{=}2$); quantitative results are in \S\ref{sec:results}.}
  \label{fig:main}
\end{figure*}

\section{Introduction}
\label{sec:intro}
% TODO: personalisation needs memory; three homes (weights, context, activations).
% NPM (2026) shows steering works as procedural memory and assumes declarative content needs text. We test it.
% Contributions: (1) a gated activation memory for user content; (2) the fixed-average account and its predictions;
% (3) pre-registered evidence on two scales with a cross-family judge panel; (4) the reference effect (dislikes flip).

\section{Method}
\label{sec:method}
% Write (shift, references: opposite / centroid / without), key (whitened pooled gist), store (delta / least squares;
% the ROME/MEMIT associative memory, in activations), read (gate as in CAST; h <- h + alpha M q at 3 late layers).

\section{Why a stored shift keeps some content and not other}
\label{sec:theory}
% In context, attention recomputes the experience's effect per question (ICV; Geva et al.).
% A stored shift is one fixed average: only question-independent content survives.
% P1 facts fail (word-to-say, late only); P2 leanings in a dose window; P3 the reference decides concept vs direction.

\section{Experimental setup}
\label{sec:setup}
% Qwen3.5-2B and 9B (fp32), thinking off; items (15 leanings, 8 dislikes, 6 one-of-many, 12 facts);
% doses, controls (random, swapped, placebo, gate forced on); pre-registration; judge + cross-family panel; bootstrap CIs.

\section{Results}
\label{sec:results}
\subsection{The memory is selective}
\subsection{Facts are stored as a late word and are not used}
\subsection{Leanings transfer only in a narrow dose window}
\subsection{What the shift subtracts decides concept vs direction}

\section{Related work}
\label{sec:related}

\section{Limitations}
\label{sec:limits}

\section{Conclusion}
\label{sec:conclusion}

\nocite{*} % TODO remove once the text cites
\bibliography{refs}

\appendix
\section{Items, prompts and judge rubrics}
\section{Judge panel and sanity check}
\section{Additional figures}
\section{Pre-registration and reproduction}

\end{document}
